% tasks T2 T3 T5 T6 T7 T8 of the Iteration 2 plan (Madhu's tasks)
% paste each section into the Overleaf template under the matching heading
% owner line is left empty on purpose, fill in who signs it
% checker for T2 T5 T6 T7 T8 is Aarav and for T3 is Margarita

%=====================================================================
\section{Notation and Definitions}
%=====================================================================

\subsection{Notation}

\begin{table}[h]
\centering
\small
\begin{tabular}{lll}
\hline
Symbol & Meaning & Size \\
\hline
$n_u$ & number of users (943 in MovieLens 100K) & number \\
$n_i$ & number of movies (1{,}682 in MovieLens 100K) & number \\
$R$ & rating matrix, one row per user, one column per movie & $n_u \times n_i$ \\
$r_{ui}$ & stars user $u$ gave to movie $i$, from 1 to 5 & number \\
$\Omega$ & set of pairs $(u,i)$ with a known rating & $|\Omega| = 100{,}000$ \\
$M$ & mask, $M_{ui}=1$ if $(u,i)\in\Omega$, else $0$ & $n_u \times n_i$ \\
$k$ & number of hidden factors per user and per movie & number \\
$U$ & user factor matrix, row $u$ is $u_u^\top$ & $n_u \times k$ \\
$u_u$ & factors of user $u$ & $k \times 1$ \\
$V$ & movie factor matrix, row $i$ is $v_i^\top$ & $n_i \times k$ \\
$v_i$ & factors of movie $i$ & $k \times 1$ \\
$\hat r_{ui}$ & guessed rating, $\hat r_{ui} = u_u^\top v_i$ & number \\
$\lambda$ & penalty weight, $\lambda > 0$ & number \\
$J(U,V)$ & objective from Question 1 & number \\
$\Omega_u$ & movies rated by user $u$, $\{i : (u,i)\in\Omega\}$ & set \\
$m_u$ & number of movies rated by user $u$, $m_u = |\Omega_u|$ & number \\
$V_u$ & rows $v_i^\top$ of $V$ for $i \in \Omega_u$ & $m_u \times k$ \\
$r_u$ & ratings $r_{ui}$ of user $u$ for $i \in \Omega_u$ & $m_u \times 1$ \\
$\Omega_i$ & users who rated movie $i$, $\{u : (u,i)\in\Omega\}$ & set \\
$m_i$ & number of users who rated movie $i$, $m_i = |\Omega_i|$ & number \\
$U_i$ & rows $u_u^\top$ of $U$ for $u \in \Omega_i$ & $m_i \times k$ \\
$r_i$ & ratings $r_{ui}$ of movie $i$ for $u \in \Omega_i$ & $m_i \times 1$ \\
$I$ & identity matrix & $k \times k$ \\
$t$ & ALS step number, $U^{t}, V^{t}$ are the values after step $t$ & number \\
$\sigma_j$ & $j$-th singular value of a matrix, $\sigma_1 \ge \sigma_2 \ge \dots \ge 0$ & number \\
$A_k$ & truncated SVD of $A$ with the top $k$ singular values & same as $A$ \\
\hline
\end{tabular}
\caption{Every symbol used in the report.}
\label{tab:notation}
\end{table}

\subsection{Definitions}

\begin{itemize}
  \item \textbf{Dot product.} For $a, b \in \mathbb{R}^k$, $a^\top b = \sum_{j=1}^{k} a_j b_j$. The guess $\hat r_{ui} = u_u^\top v_i$ is a dot product.
  \item \textbf{Euclidean norm.} For $x \in \mathbb{R}^k$, $\|x\|^2 = x^\top x = \sum_j x_j^2$.
  \item \textbf{Frobenius norm.} For a matrix $A$, $\|A\|_F^2 = \sum_{a}\sum_{b} A_{ab}^2$, the sum of the squares of all entries. So $\|U\|_F^2 = \sum_{u} \|u_u\|^2$ (sum over the rows).
  \item \textbf{Rank.} The rank of $A$ is the number of linearly independent columns of $A$. If $A = UV^\top$ with $U$ having $k$ columns, then $\operatorname{rank}(A) \le k$.
  \item \textbf{Matrix factorization.} Writing $R \approx UV^\top$ with thin $U$ and $V$ (here $k$ is much smaller than $n_u$ and $n_i$).
  \item \textbf{Symmetric matrix.} A square matrix $A$ with $A^\top = A$.
  \item \textbf{Positive semi-definite.} A symmetric $A \in \mathbb{R}^{k\times k}$ with $x^\top A x \ge 0$ for every $x \in \mathbb{R}^k$.
  \item \textbf{Positive definite.} A symmetric $A \in \mathbb{R}^{k\times k}$ with $x^\top A x > 0$ for every $x \ne 0$.
  \item \textbf{Invertible matrix.} A square $A$ is invertible if there is $A^{-1}$ with $AA^{-1} = A^{-1}A = I$. Equivalently, $Ax = 0$ only for $x = 0$ (invertible matrix theorem).
  \item \textbf{Gradient.} For $f:\mathbb{R}^k \to \mathbb{R}$, $\nabla f(x)$ is the vector of partial derivatives $\partial f / \partial x_j$.
  \item \textbf{Singular value decomposition (SVD).} Every $A \in \mathbb{R}^{m\times n}$ can be written $A = \sum_{j} \sigma_j a_j b_j^\top$ with orthonormal vectors $a_j$, orthonormal vectors $b_j$ and $\sigma_1 \ge \sigma_2 \ge \dots \ge 0$.
  \item \textbf{Truncated SVD.} $A_k = \sum_{j=1}^{k} \sigma_j a_j b_j^\top$, the SVD with only the top $k$ terms kept.
  \item \textbf{Alternating least squares (ALS).} Repeat two steps: fix $V$ and find the best $U$, then fix $U$ and find the best $V$. Each step is a least squares problem.
  \item \textbf{RMSE.} For guesses $\hat r_{ui}$ on a set of pairs $T$, \[
  \text{RMSE} = \sqrt{\frac{1}{|T|}\sum_{(u,i)\in T} (r_{ui} - \hat r_{ui})^2} .
  \]
\end{itemize}

\textbf{Data and math.} One row of $R$ is one user, one column is one movie, and one entry $r_{ui}$ is the stars user $u$ gave to movie $i$. Section Data gives the full description.

%=====================================================================
\section{Assumptions}
%=====================================================================

\begin{table}[h]
\centering
\small
\begin{tabular}{p{2.6cm}p{4.2cm}p{4.2cm}}
\hline
Assumption & What goes wrong if it fails & How we check it in Iteration 3 \\
\hline
\textbf{A1 Low rank.} Ratings are explained by a few hidden factors, so $R \approx UV^\top$ with small $k$. &
$UV^\top$ has rank at most $k$, so it cannot copy a matrix that needs more factors. The error stays high on both train and test data. &
Run ALS for $k = 1, 2, 5, 10, 20$ and plot train and test RMSE against $k$. \\[4pt]
\textbf{A2 Missing at random.} Whether a cell is in $\Omega$ does not depend on its rating. &
People mostly rate movies they chose to watch and liked, so known ratings are higher than the unknown ones would be. The guesses for empty cells come out too high. &
We cannot fully check it without the hidden values. We compare the mean rating of movies with many ratings and with few ratings. \\[4pt]
\textbf{A3 Every user and every movie has a training rating.} Each row and each column of $M$ has at least one 1. &
If $m_u = 0$ then $V_u$ is empty and the update gives $u_u = (\lambda I)^{-1} 0 = 0$, so every guess for that user is 0 stars. In \texttt{u1.base} 32 movies have no rating. &
Count the ratings per user and per movie in \texttt{u1.base}. For movies with no rating we report that the guess is not meaningful. \\[4pt]
\textbf{A4 $\lambda > 0$.} &
With $\lambda = 0$, $V_u^\top V_u$ can be singular (Question 3), so the update does not exist. 141 movies have only one rating, so $U_i^\top U_i$ has rank 1 and is singular for $k \ge 2$. &
Choose $\lambda$ on part of the training data and report test RMSE for each $\lambda$ tried. \\[4pt]
\textbf{A5 Errors are independent with the same spread.} $r_{ui} = u_u^\top v_i + \varepsilon_{ui}$. &
Squared error gives every rating the same weight. If some users are much noisier, they pull $U$ and $V$ more than the rest. &
Plot the test errors against the true star value and against the number of ratings of the user. \\[4pt]
\textbf{A6 Stars are real numbers.} We treat 1 to 5 stars as numbers on a line. &
The guess $u_u^\top v_i$ can be below 1 or above 5, which is not a real rating. &
Count the guesses outside $[1,5]$ and compare RMSE with and without clipping them to $[1,5]$. \\[4pt]
\textbf{A7 Taste does not change over time.} We ignore the time column. &
If taste changes, old ratings describe a different person, and the guess for a new rating is worse. &
Compare test RMSE for early and late ratings by timestamp. \\
\hline
\end{tabular}
\caption{Assumptions of the method.}
\label{tab:assumptions}
\end{table}

%=====================================================================
% questions 2 to 5 go after Question 1 under Iteration 2 Math Questions
%=====================================================================

\subsection*{Question 2 The update for one user}
\textit{Owner \underline{\hspace{4cm}}. Checked by \underline{\hspace{4cm}}}

\textbf{Proposition.} Fix $V$. The value of $u_u$ that makes $J$ smallest is
\[
u_u = \left(V_u^\top V_u + \lambda I\right)^{-1} V_u^\top r_u .
\]

\textbf{Step 1. Keep only the terms with $u_u$.} In $J$ from Question 1, $u_u$ appears only in the pairs $(u,i)$ with $i \in \Omega_u$ and in the penalty $\lambda\|u_u\|^2$, because $\|U\|_F^2 = \sum_{u}\|u_u\|^2$ (definition of the Frobenius norm, sum over rows). All other terms do not change when $u_u$ changes, so we call their sum $C$. So
\[
J = \sum_{i\in\Omega_u} \left(r_{ui} - v_i^\top u_u\right)^2 + \lambda \|u_u\|^2 + C ,
\]
where we used $u_u^\top v_i = v_i^\top u_u$ (a dot product does not depend on the order).

\textbf{Step 2. Write the sum as one vector.} Row $i$ of $V_u u_u$ is $v_i^\top u_u$ (definition of the matrix--vector product), and entry $i$ of $r_u$ is $r_{ui}$. So the vector $r_u - V_u u_u$ has entries $r_{ui} - v_i^\top u_u$, and the sum of their squares is its squared norm (definition of the Euclidean norm). So
\[
J = \|r_u - V_u u_u\|^2 + \lambda\, u_u^\top u_u + C .
\]

\textbf{Step 3. Expand.} We use $\|x\|^2 = x^\top x$, $(a-b)^\top(a-b) = a^\top a - a^\top b - b^\top a + b^\top b$, $(V_u u_u)^\top = u_u^\top V_u^\top$ (transpose of a product), and $r_u^\top V_u u_u = u_u^\top V_u^\top r_u$ (a number is equal to its own transpose). So
\begin{align*}
J &= r_u^\top r_u - 2\, u_u^\top V_u^\top r_u + u_u^\top V_u^\top V_u u_u + \lambda\, u_u^\top u_u + C \\
  &= r_u^\top r_u - 2\, u_u^\top V_u^\top r_u + u_u^\top \left(V_u^\top V_u + \lambda I\right) u_u + C ,
\end{align*}
where the last step uses $u_u^\top u_u = u_u^\top I u_u$.

\textbf{Step 4. Take the gradient.} Let $A = V_u^\top V_u + \lambda I$ and $b = V_u^\top r_u$. $A$ is symmetric, because $(V_u^\top V_u)^\top = V_u^\top V_u$ and $I^\top = I$. We use the rules $\nabla_x (b^\top x) = b$ and $\nabla_x (x^\top A x) = 2Ax$ for symmetric $A$. The terms $r_u^\top r_u$ and $C$ do not contain $u_u$, so their gradient is 0. So
\[
\nabla_{u_u} J = -2\, V_u^\top r_u + 2\left(V_u^\top V_u + \lambda I\right) u_u .
\]

\textbf{Step 5. Set the gradient to zero.} At a minimum the gradient is zero (first-order condition). Dividing by 2 gives the normal equations for one user
\[
\left(V_u^\top V_u + \lambda I\right) u_u = V_u^\top r_u .
\]

\textbf{Step 6. Solve.} By Question 3, $V_u^\top V_u + \lambda I$ is invertible. We multiply both sides on the left by its inverse, so
\[
u_u = \left(V_u^\top V_u + \lambda I\right)^{-1} V_u^\top r_u .
\]

\textbf{Step 7. It is a minimum, not a maximum.} The Hessian (matrix of second derivatives) is $\nabla^2_{u_u} J = 2\left(V_u^\top V_u + \lambda I\right)$, which is positive definite by Question 3. A function whose Hessian is positive definite everywhere is strictly convex, and a strictly convex function has exactly one minimum, at the point where its gradient is zero. So Step 6 gives the one and only minimizer.

\textbf{The same for one movie.} Fix $U$ instead. The same steps with the roles of users and movies swapped give
\[
v_i = \left(U_i^\top U_i + \lambda I\right)^{-1} U_i^\top r_i .
\]

\textbf{Conclusion.} With $V$ fixed, the best factors for user $u$ come from one $k \times k$ linear system. In the code we solve this system with \texttt{np.linalg.solve(left, right)}, where \texttt{left} $= V_u^\top V_u + \lambda I$ and \texttt{right} $= V_u^\top r_u$ (\texttt{als} in \texttt{src/simulate.py}). We never compute the inverse itself.

\subsection*{Question 3 The matrix is positive definite}
\textit{Owner \underline{\hspace{4cm}}. Checked by \underline{\hspace{4cm}}}

\textbf{Proposition.} For $\lambda > 0$, $A = V_u^\top V_u + \lambda I$ is positive definite, so it is invertible.

\textbf{Step 1. $A$ is symmetric.} $(V_u^\top V_u)^\top = V_u^\top (V_u^\top)^\top = V_u^\top V_u$ (transpose of a product reverses the order, and $(X^\top)^\top = X$), and $I^\top = I$. A sum of symmetric matrices is symmetric. So $A^\top = A$.

\textbf{Step 2. Compute $x^\top A x$.} Take any $x \in \mathbb{R}^k$ with $x \ne 0$. By the distributive law,
\[
x^\top A x = x^\top V_u^\top V_u x + \lambda\, x^\top x = (V_u x)^\top (V_u x) + \lambda\, x^\top x = \|V_u x\|^2 + \lambda \|x\|^2 ,
\]
using $(V_u x)^\top = x^\top V_u^\top$ and the definition of the Euclidean norm.

\textbf{Step 3. It is positive.} $\|V_u x\|^2 \ge 0$, because it is a sum of squares. $\|x\|^2 > 0$, because $x \ne 0$ has at least one entry that is not zero. With $\lambda > 0$ we get $\lambda\|x\|^2 > 0$. So $x^\top A x \ge 0 + \lambda \|x\|^2 > 0$, and $A$ is positive definite (definition).

\textbf{Step 4. Positive definite gives invertible.} Suppose $Ax = 0$ for some $x \ne 0$. Then $x^\top A x = x^\top 0 = 0$, which contradicts Step 3. So $Ax = 0$ only for $x = 0$, and the square matrix $A$ is invertible (invertible matrix theorem).

\textbf{Why $\lambda$ is needed.} Without $\lambda$, $V_u^\top V_u$ is only positive semi-definite. Its rank is at most $m_u$, because $V_u$ has $m_u$ rows. If a user or a movie has fewer than $k$ ratings, the $k \times k$ matrix has rank below $k$ and is singular. In MovieLens 100K, 141 movies have only one rating, so for them $U_i^\top U_i$ has rank 1 and cannot be inverted for $k \ge 2$. The proof above works even when $m_u = 0$, because then $A = \lambda I$.

\textbf{Conclusion.} For every $\lambda > 0$ the system in Question 2 has exactly one solution, for every user and every movie, no matter how few ratings they have.

\subsection*{Question 4 An ALS step never increases $J$}
\textit{Owner \underline{\hspace{4cm}}. Checked by \underline{\hspace{4cm}}}

\textbf{Proposition.} Let $U^{t+1}$ be all user updates of Question 2 with $V^{t}$ fixed, and $V^{t+1}$ be all movie updates with $U^{t+1}$ fixed. Then
\[
J(U^{t+1}, V^{t+1}) \le J(U^{t+1}, V^{t}) \le J(U^{t}, V^{t}).
\]

\textbf{Step 1. $J$ splits into one part per user.} With $V$ fixed, every term of $J$ contains at most one row $u_u$ of $U$. We group the terms by user, using $\Omega = \bigcup_u \{(u,i) : i\in\Omega_u\}$ and $\|U\|_F^2 = \sum_u \|u_u\|^2$:
\[
J(U, V) = \sum_{u=1}^{n_u} \underbrace{\left( \sum_{i\in\Omega_u} \left(r_{ui} - u_u^\top v_i\right)^2 + \lambda \|u_u\|^2 \right)}_{J_u(u_u)} + \lambda \|V\|_F^2 .
\]
So $J_u$ depends only on $u_u$, and the last term does not depend on $U$ at all.

\textbf{Step 2. The user step finds the best $U$.} By Questions 2 and 3, $u_u^{t+1}$ is the one minimizer of $J_u$ with $V = V^t$. So $J_u(u_u^{t+1}) \le J_u(u_u)$ for every $u_u$. Adding these inequalities over all users (a sum of smaller numbers is smaller), and adding the same $\lambda\|V^t\|_F^2$ to both sides,
\[
J(U^{t+1}, V^{t}) \le J(U, V^{t}) \quad \text{for every } U .
\]
We choose $U = U^{t}$, which gives $J(U^{t+1}, V^{t}) \le J(U^{t}, V^{t})$.

\textbf{Step 3. The movie step finds the best $V$.} With $U = U^{t+1}$ fixed, $J$ splits into one part per movie in the same way, and $v_i^{t+1}$ is the one minimizer of each part (Question 2 for one movie, Question 3 with $U_i$). The same argument gives $J(U^{t+1}, V^{t+1}) \le J(U^{t+1}, V^{t})$.

\textbf{Step 4. Put the two steps together.} By Steps 2 and 3 and the transitive rule ($a \le b$ and $b \le c$ give $a \le c$),
\[
J(U^{t+1}, V^{t+1}) \le J(U^{t+1}, V^{t}) \le J(U^{t}, V^{t}).
\]

\textbf{Step 5. The values of $J$ converge.} The numbers $J(U^t, V^t)$ never go up (Step 4) and never go below 0 (Question 1, $J$ is a sum of squares with $\lambda > 0$). A sequence that never goes up and is bounded below converges (monotone convergence theorem). So $J$ settles at some value.

\textbf{What this does not say.} It does not say that ALS finds the smallest possible $J$. $J$ is not convex in $U$ and $V$ together, because of the product $u_u^\top v_i$. So ALS can stop at a value that depends on the random start.

\textbf{Conclusion.} Every ALS step makes $J$ smaller or keeps it the same. The simulation checks this: the largest change of $J$ between two steps is $-2.4\times10^{-5}$ (Table~\ref{tab:sim}).

\subsection*{Question 5 Truncated SVD with missing ratings}
\textit{Owner \underline{\hspace{4cm}}. Checked by \underline{\hspace{4cm}}}

\textbf{Claim.} The truncated SVD is the best rank-$k$ guess when every entry of the matrix is known, but not when ratings are missing.

\textbf{Step 1. Eckart--Young for a full matrix.} The Eckart--Young theorem says: for any $A \in \mathbb{R}^{m\times n}$ with SVD $A = \sum_j \sigma_j a_j b_j^\top$,
\[
\min_{\operatorname{rank}(B)\le k} \|A - B\|_F^2 = \|A - A_k\|_F^2 = \sum_{j>k} \sigma_j^2 .
\]
So the truncated SVD $A_k$ is the best rank-$k$ matrix, and its error is the sum of the squares of the dropped singular values.

\textbf{Step 2. When it solves our problem.} Suppose every rating is known, so $M$ is all ones, and $\lambda = 0$. Then $J(U,V) = \|R - UV^\top\|_F^2$ (Question 1, matrix form). Every $UV^\top$ with $k$ columns has rank at most $k$, and every matrix $B$ of rank at most $k$ can be written as $UV^\top$ (for example $U = [\sigma_1 a_1, \dots, \sigma_k a_k]$, $V = [b_1, \dots, b_k]$ from the SVD of $B$). So minimizing $J$ is the same as Step 1, and the answer is $UV^\top = R_k$. In this case one SVD gives the best answer directly, without iteration.

\textbf{Step 3. Why it breaks with missing ratings.}
\begin{itemize}
  \item \textbf{The SVD needs every entry.} The SVD is defined only for a full matrix. 93.7\% of $R$ is unknown (Section Data), so we must first fill the empty cells with some number, for example 0 or the mean 3.53.
  \item \textbf{It then minimizes the wrong objective.} After filling, the SVD minimizes $\|R_{\text{filled}} - B\|_F^2$, a sum over all $1{,}586{,}126$ cells. Only $100{,}000$ of them are real ratings, and the other $1{,}486{,}126$ are values we made up. So the best $B$ mostly copies the made-up values, not the ratings. Our objective is $\|M \odot (R - B)\|_F^2$, which counts only the cells in $\Omega$, and that is a different problem.
  \item \textbf{The proof of Eckart--Young does not carry over.} The proof uses that the Frobenius norm does not change under multiplication by an orthogonal matrix $Q$, $\|QA\|_F = \|A\|_F$. With the mask, $M \odot (QA) \ne Q(M \odot A)$ in general, so this step fails, and there is no formula like $\sum_{j>k}\sigma_j^2$ for the masked problem.
\end{itemize}

\textbf{Step 4. What we do instead.} Because there is no direct formula, we minimize the masked objective $J$ step by step with ALS. Each step uses only known ratings (Question 2) and never makes $J$ larger (Question 4).

\textbf{Conclusion.} The truncated SVD is optimal only for the full-matrix problem. With 93.7\% of ratings missing, it must fit made-up values, so it is not optimal for our problem. The simulation shows this: the filled SVD has RMSE 1.394 on hidden cells, and ALS has 0.147 (Table~\ref{tab:sim}).
